\chapter{双主线工程图纸检索方法}

在系统总体结构确定之后，本文进一步从方法层面对工程图纸检索模型进行展开。本章重点回答两个问题：一是如何围绕工程图纸“结构主导、文本辅助”的特点构建更适合的检索路线；二是如何在统一实验框架下对显式清洗与无清洗主体增强两种思路进行可比较的设计。为此，本文将原有串行堆叠式流程重构为双主线方法体系，并在两条主线之上分别考察结构增强与 OCR 辅助的作用。

\section{方法设计思路}

针对工程图纸检索任务中存在的背景干扰、局部结构相似、标题栏文本信息分散等问题，本文不再沿用“清洗、局部增强、结构重排、OCR 融合”单链条串行堆叠的叙述方式，而是将整个方法重构为两条主线：
\begin{enumerate}
  \item \textbf{YOLO 清洗路线}：先通过目标检测清除无关区域，再在清洗结果上进行视觉建模、结构增强和 OCR 辅助检索；
  \item \textbf{无清洗注意力路线}：不显式执行 YOLO 清洗，而是直接利用掩码引导的局部注意力式建模突出主体结构，再叠加结构增强与 OCR 辅助检索。
\end{enumerate}

这样的设计可以更清楚地回答三个研究问题：第一，YOLO 清洗是否在新数据集上仍然有效；第二，不依赖显式清洗的注意力式建模是否能够取得更优结果；第三，两条路线在引入结构增强和 OCR 辅助后，哪一条路线的收益更稳定。

从算法组织角度看，本文方法包含“候选召回”和“候选精排”两个阶段。候选召回阶段主要依赖 CLIP 视觉向量，以保证系统能够在较大图库中快速获得相似候选；候选精排阶段再根据具体模式引入结构相似度和 OCR 字段相似度，对前列结果进行细粒度排序修正。这样的两阶段设计既保留了向量检索的效率优势，也为工程图纸中重要的局部结构和标题栏信息预留了融合空间。

\section{基础视觉检索框架}

本文首先构建基于 CLIP 视觉编码器的工程图纸检索基线。对于输入图纸 $I$，视觉编码器输出全局特征 $\bm{f}(I)$，归一化后得到检索向量：
\begin{equation}
  \bm{z}(I)=\frac{\bm{f}(I)}{\|\bm{f}(I)\|_2}.
  \label{eq:baseline-embedding-new}
\end{equation}

对于查询图纸 $I_q$ 与图库图纸 $I_i$，语义相似度定义为
\begin{equation}
  s_{\text{sem}}(I_q,I_i)=\bm{z}(I_q)^\top \bm{z}(I_i).
  \label{eq:semantic-score-new}
\end{equation}

该基线能够提供稳定的全局视觉召回能力，但对工程图纸中的细粒度结构布局、标题栏字段信息和局部主体区域关注不足，因此需要在其基础上进一步增强。

\section{YOLO 清洗路线}

\subsection{检测清洗与保留掩码}

在 YOLO 清洗路线中，首先利用目标检测模型识别图纸中的无关干扰区域，例如截图边框、注释噪声和非主体块区域。对于检测到的每个候选框，采用填白策略将对应区域置为背景，并同步生成保留掩码。设原图为 $I$，清洗后的图像为 $\widetilde{I}$，二值保留掩码为 $M$，则
\begin{equation}
  M(x,y)=
  \begin{cases}
    1, & (x,y)\ \text{属于保留区域},\\
    0, & (x,y)\ \text{属于清洗区域}.
  \end{cases}
\end{equation}

这里的掩码并不直接替代原始视觉编码，而是作为后续局部结构增强的先验信息。实验部分将专门验证：在新数据集已经完成黑底和黑边统一处理后，单独执行 YOLO 清洗是否仍然带来稳定收益。

\subsection{清洗路线的局部视觉增强}

在完成清洗后，本文进一步在视觉编码阶段使用掩码引导的 patch 聚合。设视觉编码器输出的 patch 特征为 $\{\bm{h}_p\}_{p=1}^{P}$，将图像级掩码映射到 patch 空间，得到 patch 权重 $m_p$，则局部增强后的特征写为
\begin{equation}
  w_p=\frac{m_p+\epsilon}{\sum_{j=1}^{P}(m_j+\epsilon)},
\end{equation}
\begin{equation}
  \bm{z}_{\text{mask}}=\sum_{p=1}^{P}w_p\bm{h}_p.
  \label{eq:masked-pooling-new}
\end{equation}

其中 $\epsilon$ 为平滑项。该机制可降低大面积空白区域与边框区域对全局特征的干扰，使模型更加关注轮廓、剖视、连接关系等主体结构信息。

\section{无清洗注意力路线}

\subsection{无清洗条件下的主体增强}

与 YOLO 清洗路线不同，无清洗注意力路线不显式对图纸做目标检测清除，而是直接从原始图纸出发，通过掩码引导的局部聚合机制突出主体结构区域。其核心思想是：即使不执行显式清洗，只要视觉编码阶段能够更聚焦主体区域，仍然可能取得优于清洗路线的检索性能。

在当前实现中，该路线采用 \textbf{mask-guided pooling / attention-like} 机制，即利用图像前景区域和局部结构分布构造轻量化的注意力式聚合，而非单独训练新的端到端空间注意力模型。也就是说，本文实现并验证的是“无清洗注意力式视觉增强路线”，而非严格意义上重新训练得到的 \texttt{spatial\_attention} 新模型。

\subsection{当前实现边界}

需要指出的是，当前实现支持将掩码引导聚合与局部 patch 表征结合，从而形成一条可运行的“无清洗注意力路线”。该方案能够作为论文中的对比路线，并用于回答“是否必须依赖显式清洗”的研究问题。

但如果要构建严格意义上的独立 \texttt{spatial\_attention} 模型，还需要单独训练与其匹配的权重文件。本文当前实验并未完成这一额外训练，因此在论文表述中应明确：本文验证的是“无清洗注意力式增强思路”，而非“重新训练的独立空间注意力模型”。

\section{结构相似性增强与重排序}

工程图纸检索与自然图像检索的重要区别在于：许多样本的局部纹理差异较小，但整体几何布局、投影轮廓和前景分布具有较强判别性。为此，本文设计了轻量级结构描述符 $\bm{g}(I)$，包括网格占用分布、水平投影、垂直投影和前景比例等信息。对于查询图纸和候选图纸，其结构相似度定义为
\begin{equation}
  s_{\text{str}}(I_q,I_i)=\operatorname{cos}\bigl(\bm{g}(I_q),\bm{g}(I_i)\bigr).
  \label{eq:structure-score-new}
\end{equation}

在候选重排序阶段，将视觉语义相似度与结构相似度融合，得到最终的结构增强评分：
\begin{equation}
  s_{\text{rerank}}=\lambda s_{\text{sem}}+(1-\lambda)s_{\text{str}},
  \label{eq:rerank-score-new}
\end{equation}
其中 $\lambda$ 为融合权重。

该设计有两个优点：一方面不会显著增加大规模图库召回成本，因为结构相似度主要作用于候选精排阶段；另一方面可以有效提升前列结果的结构一致性，尤其适用于零件图、装配图和工序图之间存在局部轮廓相似但布局关系不同的场景。

在实际计算中，结构描述符并不追求完整表达所有 CAD 几何细节，而是强调低成本、可复现和对检索排序有效。网格占用分布用于描述主体结构在图面中的空间位置，水平投影和垂直投影用于刻画轮廓延展方向，前景比例用于区分结构密集图纸与结构稀疏图纸。上述信息虽然简单，但能够补足 CLIP 全局向量在二维工程图纸布局表达上的不足。

\section{OCR 文字辅助检索}

\subsection{标题栏字段建模}

在视觉相似度与结构相似度之外，工程图纸还包含标题栏、图号、零件名称、比例和材料等强语义信息。本文没有将 OCR 输出简单视为一段普通文本，而是进一步面向标题栏抽取以下关键字段：
\begin{enumerate}
  \item 图号 $drawing\_no$；
  \item 零件名称 $part\_name$；
  \item 材料 $material$；
  \item 比例 $scale$；
  \item 数量 $quantity$。
\end{enumerate}

在此基础上，先计算通用 OCR 文本相似度 $s_{\text{ocr}}$，再计算字段级匹配得分：
\begin{equation}
  s_{\text{field}}=\alpha_1 s_{\text{id}}+\alpha_2 s_{\text{name}}+\alpha_3 s_{\text{mat}}+\alpha_4 s_{\text{scale}},
  \label{eq:field-score-new}
\end{equation}
其中 $s_{\text{id}}$、$s_{\text{name}}$、$s_{\text{mat}}$ 和 $s_{\text{scale}}$ 分别表示图号、零件名称、材料和比例的匹配得分，$\alpha_k$ 为对应权重。由于图号和零件名称通常具有更强区分度，因此其权重高于材料和比例。

\subsection{多模态融合策略}

在重排序基础上，定义 OCR 辅助后的最终得分为
\begin{equation}
  s_{\text{final}} = s_{\text{rerank}} + \beta_1 s_{\text{ocr}} + \beta_2 s_{\text{field}},
  \label{eq:final-score-new}
\end{equation}
其中 $\beta_1$ 和 $\beta_2$ 分别控制通用文本相似度和字段级匹配得分的影响。

为了避免 OCR 噪声破坏已有视觉排序，本文采用“候选重排而非全图库覆盖”的策略，仅在前列候选中启用 OCR 增强，并设置字段质量门控机制。由此，OCR 并不作为独立检索主干，而是作为视觉检索的辅助微调模块。

\section{双主线实验矩阵}

基于上述设计，本文将正式实验组织为 8 个模式：
\begin{enumerate}
  \item \texttt{baseline}；
  \item \texttt{cleaning\_only}；
  \item \texttt{masked\_pooling}；
  \item \texttt{full\_model}；
  \item \texttt{multimodal\_text}；
  \item \texttt{attention\_visual}；
  \item \texttt{attention\_structure}；
  \item \texttt{attention\_multimodal}。
\end{enumerate}

其中前 5 个模式构成 YOLO 清洗路线，后 4 个模式构成无清洗注意力路线，\texttt{baseline} 作为两条路线的共同参考点。这样的实验组织方式可以从方法设计上清楚划分“显式清洗”和“无清洗注意力式建模”两种思路，并进一步比较结构增强和 OCR 辅助在两条路线中的作用差异。

\section{方法流程总结}

综合以上设计，本文方法的完整流程可以概括为以下步骤：
\begin{enumerate}
  \item 输入查询图纸，并根据所选模式决定是否执行 YOLO 清洗；
  \item 提取查询图纸的 CLIP 视觉特征，在需要时执行掩码引导聚合或无清洗注意力式主体增强；
  \item 在 Qdrant 中利用视觉向量完成候选召回，获得初始 Top-K 候选集合；
  \item 对候选集合计算结构描述子相似度，并与视觉相似度融合得到结构重排分数；
  \item 当模式启用 OCR 时，抽取标题栏字段并计算字段级文本匹配分数；
  \item 综合视觉、结构和 OCR 分数得到最终排序结果，并输出相似图纸及辅助分析信息。
\end{enumerate}

该流程体现了本文“视觉为主、结构增强、文本辅助”的方法原则。其中视觉特征负责提供稳定召回，结构相似性负责改善前列候选的几何一致性，OCR 字段级信息则作为可扩展的辅助信号参与排序修正。通过双主线实验设计，本文能够进一步分析不同模块在真实工程图纸数据集上的实际贡献。

\section{可解释性分析}

为了辅助分析误检原因和局部关注区域，本文在检索结果分析中引入 Grad-CAM 可解释性工具。具体做法是对查询图与候选图的相似度响应执行梯度回传，在 patch 级特征上生成热力分布，从而观察模型究竟关注的是主体轮廓、剖视区域，还是无关边框与空白背景。

Grad-CAM 本身并不参与检索得分计算，但可为“局部注意力式建模是否真正聚焦主体结构”提供定性证据，也为后续误检分析与模型改进提供依据。

\section{本章小结}

本章围绕工程图纸检索任务给出了双主线方法设计。首先构建基于 CLIP 的基础视觉检索框架，然后分别提出 YOLO 清洗路线与无清洗注意力式视觉增强路线，并在此基础上引入结构重排序、OCR 字段级辅助和可解释性分析机制。通过这种组织方式，本文不仅完成了方法层面的模块化设计，也为下一章的正式实验提供了清晰的对比对象与评测维度。
